% F13 -- Where the gain lands once the group physically routes the sequence.
\begin{figure}[t]
\centering
\replegend{\swatch{clbase}~persistence baseline \qquad
           \swatch{clmodel}~model, group routed to its own adapter}
\plotlink{\nbllm}{\fitwidth{%
\begin{tikzpicture}
\begin{axis}[repbar, width=\linewidth, height=5.8cm,
  symbolic x coords={G0,G1,G2,G3,MEAN}, xtick=data,
  xticklabels={G0 homebody,G1 steady,G2 evening,G3 morning,mean of the four},
  /pgf/bar width=15pt,
  ymin=0, ymax=80, ytick={0,10,...,60},
  ylabel={top-1 accuracy (\%)},
  x tick label style={font=\small},
]
\addplot[fill=clbase,draw=none]  coordinates {(G0,56.0)(G1,42.0)(G2,19.7)(G3,17.8)(MEAN,33.9)};
\addplot[fill=clmodel,draw=none] coordinates {(G0,56.8)(G1,44.3)(G2,25.1)(G3,24.2)(MEAN,37.6)};
\node[font=\small,text=clgain,anchor=south] at (axis cs:G0,68) {$+0.8$};
\node[font=\small,text=clgain,anchor=south] at (axis cs:G1,56) {$+2.3$};
\node[font=\small,text=clgain,anchor=south] at (axis cs:G2,37) {$+5.4$};
\node[font=\small,text=clgain,anchor=south] at (axis cs:G3,36) {$+6.4$};
\node[font=\small,text=clgain,anchor=south] at (axis cs:MEAN,49.5) {$+3.7$};
\end{axis}
\end{tikzpicture}}}
\caption{Per-group gain after routing the group to its own adapter}
\label{fig:group-gain}
\figtext{%
The gain ranks in exactly the reverse order of the baseline. One pair of bars
per group, plus the unweighted mean of the four, with the persistence baseline
in grey-blue and the model in dark blue, scored on the same positions of the
same 100 held-out users over the full sequence. Nothing separates this run
from the one that scored $-1.8\pt$ except which weights were allowed to answer
(Figure~\ref{fig:adapter-routing}). G0, whom repetition already predicts
56.0\,\% of the time, gains $+0.8\pt$ and has almost nothing left to win; G3,
whom repetition predicts less than one time in five, gains $+6.4\pt$. The
ordering is exact, with no exception among the four, which is the claim; over
four points a correlation coefficient would add nothing to it. Read against
Chapter~\ref{ch:user} this is the
same finding arriving from the other direction: difficulty lives in the user,
and a model earns its cost precisely where the trivial rule fails. Two
qualifications belong with the headline. The $+3.7\pt$ is the unweighted mean
over four profiles counted once each; weighted by how many users each actually
holds it is $+3.3\pt$, which is the figure an operator would experience,
because 28\,\% of the population sits in the group that needed no model. And
the whole mechanism depends on the group being detectable, which
Figure~\ref{fig:group-difficulty}(b) shows collapsing below roughly 60\,\% of
a day of history. A third qualification, that four adapters are also four
times the trainable capacity and the control was never run, is in
Section~\ref{sec:retro-groups}.}
\figsource{\nbllm}{train\_cellid\_llm.ipynb}{the per-group adapter-routing run}
\end{figure}
